\documentclass[a4paper,11pt]{article}
\usepackage[utf8]{inputenc}
\usepackage[T1]{fontenc}
\usepackage[french]{babel}
\usepackage{geometry}
\geometry{margin=2.5cm}
\usepackage{hyperref}
\usepackage{booktabs}
\usepackage{longtable}
\usepackage{graphicx}
\usepackage{amsmath}
\usepackage{xcolor}

\title{Rapport d'experimentation\\
\large Etude systematique de la planarite des solutions CSP\\
pour les benzenoides non-benzenoides}
\author{Toky}
\date{Avril 2026}

\begin{document}
\maketitle

\begin{abstract}
Ce rapport presente les resultats de l'exploration experimentale systematique
du solveur CSP pour les substitutions pentagone/heptagone dans les benzenoides.
L'objectif est d'etudier les proprietes des solutions CSP --- nombre de solutions,
distribution des tailles de cycles, planarite --- en fonction de la taille, de la forme,
de la symetrie et des motifs des benzenoides d'entree. Les benzenoides testes proviennent
de la base BenzDB (exhaustive pour h$\leq$9) et de generateurs parametriques.
\end{abstract}

\tableofcontents
\newpage

%% ===================================================================
\section{Introduction}

\subsection{Cadre experimental}

\textbf{Outils :}
\begin{itemize}
    \item Solveur CSP : PyCSP3 + ACE (contraintes C1--C4, table de voisinage T(n))
    \item Table T(n) : T(5)=92, T(6)=418, T(7)=699 entrees (1209 total)
    \item Validation : reconstruction 3D + optimisation xTB (GFN2-xTB)
    \item Seuil de planarite : $\theta_{\max} \leq 10\degree$
\end{itemize}

\textbf{Metriques par benzenoide :}
\begin{itemize}
    \item h (nb hexagones), $|V|$, $|E|$, degres du dual, pattern $\sigma(v)$, b(v)
    \item Nb hexagones geles/libres, nb hexagones internes
    \item Nb generateurs d'automorphismes, diametre du dual
\end{itemize}

\textbf{Metriques par resolution CSP :}
\begin{itemize}
    \item Nb solutions (avec/sans freeze), temps de resolution
    \item Distribution des tailles (pentagones/hexagones/heptagones)
\end{itemize}

\textbf{Metriques par solution validee :}
\begin{itemize}
    \item Angle max, RMSD plan, hauteur, statut plan/non plan
\end{itemize}

\subsection{Sources de donnees}
\begin{itemize}
    \item BenzDB (these Varet 2022) : tous les benzenoides $h \leq 9$ (8392 molecules)
    \item Generateur parametrique : formes lineaires, rectangulaires, coronenoides, etc.
    \item Reference : these Varet (2022), BenzAI
\end{itemize}

%% ===================================================================
\section{Serie 0 : Validation sur BenzDB (h$\leq$9)}

\textit{A completer avec les resultats.}

\subsection{S0a : h=3..5 (26 benzenoides)}
%% Tableau des resultats

\subsection{S0b : h=6..7 (412 benzenoides)}
%% Tableau des resultats

\subsection{S0c : h=8..9 (7946 benzenoides)}
%% Tableau des resultats

%% ===================================================================
\section{Series thematiques}

\subsection{Catacondenses}
%% Resultats

\subsection{Rectangulaires}
%% Resultats

\subsection{Coronenoides}
%% Resultats

\subsection{Symetries}
%% Resultats

\subsection{Coronoides (avec trous)}
%% Resultats

\subsection{Motifs de bordure}
%% Resultats

\subsection{Parametre d'irregularite}
%% Resultats

%% ===================================================================
\section{Observations transversales}

\textit{A completer : comparaisons inter-categories, tendances generales.}

%% ===================================================================
\section{Resultats remarquables}

\textit{A completer : cas surprenants, patterns inattendus.}

%% ===================================================================
\section{Pistes ouvertes}

\textit{A completer : questions soulevees par les resultats.}

\end{document}
